\paragraph{Repeated Controlled Operations}\label{ex:repeated-controlled-operations}

In the previous sections, we saw that if the target qubit is an eigenstate of a unitary operator $U$:
\begin{equation*}
    U \ket{v} = e^{i \phi} \ket{v}
\end{equation*}
Then a controlled application of $U$ transfers the eigenvalue phase $e^{i \phi}$ onto the control qubit through the \emph{phase kickback} effect. Now, after having seen how to measure the kicked-back phase, we can ask ourselves \emph{\textbf{what happens if we apply the controlled operation multiple times?}} The answer is really simple: \textbf{each application contributes the same phase, so the phases accumulate}. This idea is fundamental because it is exactly what allows algorithms such as Quantum Phase Estimation (QPE) to determine an unknown phase with very high precision.

\highspace
\begin{flushleft}
    \textcolor{Green3}{\faIcon{redo} \textbf{Applying a unitary twice}}
\end{flushleft}
Let's start with the application of a controlled unitary $U$:
\begin{equation*}
    U\ket{v} = e^{i \phi} \ket{v}
\end{equation*}
Where:
\begin{itemize}
    \item \ket{v} is an eigenstate of $U$;
    \item $e^{i \phi}$ is the corresponding eigenvalue, which is a phase factor.
\end{itemize}
Now apply the same operator a \textbf{second time}. We obtain:
\begin{equation*}
    U^2 \ket{v} = U \left(U \ket{v}\right)
\end{equation*}
Substitute the eigenvalue equation:
\begin{equation*}
    U^2 \ket{v} = U \left(e^{i \phi} \ket{v}\right)
\end{equation*}
Since $e^{i \phi}$ is just a complex scalar, it can be factored outside the operator:
\begin{equation*}
    U^2 \ket{v} = e^{i \phi} U \ket{v}
\end{equation*}
Apply the eigenvalue equation once more:
\begin{equation*}
    U^2 \ket{v} = e^{i \phi} \left(e^{i \phi} \ket{v}\right) = e^{2 i \phi} \ket{v}
\end{equation*}
\textcolor{Green3}{\faIcon{redo} \textbf{Applying a unitary $n$ times}}. After applying the unitary $n$ times:
\begin{equation*}
    U^n = \underbrace{U \, U \, \dots \, U}_{n \text{ times}}
\end{equation*}
The eigenstate satisfies the following equation:
\begin{equation}
    U^n \ket{v} = e^{i n \phi} \ket{v}
\end{equation}
The key observation is that \textbf{the state never stops being an eigenstate}. Multiplying a vector by a scalar does \textbf{not} change its direction in Hilbert space (it only changes its phase). Therefore, after every application, the state is still proportional to \ket{v}, so applying $U$ again produces exactly the same eigenvalue. This process repeats indefinitely.

\highspace
\begin{flushleft}
    \textcolor{Green3}{\faIcon{book} \textbf{Controlled powers of a unitary}}
\end{flushleft}
Now recall the phase kickback circuit. Previously we saw that if the target qubit is an eigenstate of $U$, then a controlled application of $U$ transfers the eigenvalue phase $e^{i \phi}$ onto the control qubit. Now, suppose instead that we use a controlled version of $U^n$. The control qubit starts in:
\begin{equation*}
    \ket{+} = \frac{1}{\sqrt{2}} \left(\ket{0} + \ket{1}\right)
\end{equation*}
The target is still the eigenstate \ket{v}. Initially, the state of the two qubits is:
\begin{equation*}
    \ket{\psi} = \frac{1}{\sqrt{2}} \left(\ket{0}\ket{v} + \ket{1}\ket{v}\right) = \frac{1}{\sqrt{2}} \left(\ket{0} + \ket{1}\right) \otimes \ket{v}
\end{equation*}
Applying the controlled operator gives:
\begin{equation*}
    CU^{n}\ket{\psi} = \frac{\ket{0}\ket{v} + \ket{1} U^{n} \ket{v}}{\sqrt{2}}
\end{equation*}
Using the eigenvalue equation for $U^n$:
\begin{equation*}
    U^{n} \ket{v} = e^{i n \phi} \ket{v}
\end{equation*}
We can substitute this into the previous equation:
\begin{equation}
    CU^{n}\ket{\psi} = \left(\frac{\ket{0} + e^{i n \phi} \ket{1}}{\sqrt{2}}\right) \otimes \ket{v}
\end{equation}
Notice that the kicked-back phase is now $e^{i n \phi}$, instead of $e^{i \phi}$. This is exactly what we expected: \hl{the phase accumulates with each application of the unitary operator.}

\highspace
\begin{flushleft}
    \textcolor{Green3}{\faIcon{check-circle} \textbf{Why is this useful?}}
\end{flushleft}
At first glance, multiplying a phase might not seem very exciting. However, it becomes extremely powerful when we want to \textbf{estimate an unknown phase}.

\highspace
Suppose:
\begin{equation*}
    U\ket{v} = e^{i \phi} \ket{v}
\end{equation*}
Where $\phi$ \textbf{is completely unknown}. After \emph{phase kickback} we obtain the state:
\begin{equation*}
    \ket{\psi} = \frac{1}{\sqrt{2}} \left(\ket{0} + e^{i \phi} \ket{1}\right)
\end{equation*}
Great, but \emph{\textbf{can we read $\phi$ by measuring the qubit?}} The answer is no. If we measure in the computational basis, we will always get either $\ket{0}$ or $\ket{1}$ with equal probability (50\%). The probabilities do \textbf{not} depend on $\phi$. So one controlled operation is \textbf{not enough} to estimate the phase. Also, if we apply the controlled operation multiple times, we can accumulate the phase:
\begin{equation*}
    CU^{n}\ket{\psi} = \frac{1}{\sqrt{2}} \left(\ket{0} + e^{i n \phi} \ket{1}\right)
\end{equation*}
And the probabilities of measuring $\ket{0}$ or $\ket{1}$ are still 50\% each. So even if we apply the controlled operation multiple times, we still cannot estimate the phase.

\highspace
However, if we apply \textbf{different powers of the unitary} ($U^1, U^2, U^4, \dots$), we can \textbf{accumulate different phases on different qubits} (i.e., the phase is accumulated on each control qubit). This is exactly what \hl{Quantum Phase Estimation (QPE)} does: it \hl{uses multiple control qubits, each controlling a different power of the unitary, to accumulate different phases on each control qubit.} Then, by measuring all the control qubits, we can estimate the unknown phase $\phi$ with very high precision.